% WORKING SAME-FAMILY APPENDIX, 25 October 2026.
% Current batch: complete A16 + actual typed-seven consumer only.
% Unintegrated, uncompiled, and awaiting different-author/root review.
% Source: compact18_a16_actual_typed_seven_endpoint_defect_coupling_exclusion_candidate_20261022.md.
% No scientific import, enumeration, test, Lean or TeX job.
\section{The other exact-dimension-eighteen consumers}
\label{app:otherexact18}
The necessary forest cover is proved before any row is excluded.
This appendix supplies its component-specific consumers, using the same
original farthest vertex throughout.  A typed exception includes its
actual isolated opposite wall; a bare old seven bond is insufficient.

\subsection{\(A_{16}\) with either actual typed-seven defect}
\label{app:A16typed18}
Suppose the entire incident forest is \(A_{16}+Z\), where
\(Z=N_7\) means old \(A_2\) with its new isolated seven opposite,
and \(Z=O_7\) means old \(I_2(7)\) with its new isolated three
opposite.  There is no arithmeticity hypothesis.

Use zero-sum \(\mathbb R^{17}\), with actual unit normals and duals
\[
 n_i=(e_i-e_{i+1})/\sqrt2,\qquad
 \Omega_i=\sqrt2(P_i-iP_{17}/17),\qquad
 D_i=2i(17-i)/17,\quad C_{ij}=2i(17-j)/17\quad(i\le j).
\]
Write the original interior point as \(z=E(p-x-y)\), with
\(x\) in this \(A_{16}\) space and \(y\) in the orthogonal
rank-two space.  The previously proved typed-defect floor gives
\[
 e_A+e_Z=\norm x^2+\norm y^2<1,\qquad e_Z>.09,
 \qquad e_A<.91.                                      \tag{T1}
\]
All auxiliary points below are compared with this original maximum.

Put \(\phi=(1+\sqrt5)/2\), \(\theta=2\cos(2\pi/7)\),
\(g_7=\cos(\pi/7)\), and
\[
 h_N=(\theta-1)/3,\qquad h_O=(\theta-1)/(2-\theta),
 \qquad k=1/h_O=\theta^2+2\theta-1.
\]
The actual isolated defect wall has time \(\sqrt{h_Z}\) and
spatial support only in \(Z\).  In the \(N_7\) case it is
\(c_Z=\sqrt{h_N}p-g_7\omega_a^Z\), for an old \(A_2\) dual.
The cubic \(f(t)=t^3+t^2-2t-1\), increasing on \([1.24,1.25]\),
has signs
\[
 f(1.2469)=-410632291/10^{12}<0,\qquad
 f(1.247)=105223/10^9>0.
\]
Consequently
\[
 1.2469<\theta<1.247,\quad 0<h_N<1/12,\qquad
 3.0485<k<3.0491<61/20.                              \tag{T2}
\]
The identity \((\theta-1)k=2-\theta\) follows by the same cubic.
These are real native-angle facts, not field restrictions on the polytope.

\subsubsection{The central window at its actual energy budget}
Call \(5\le i\le12\) weak if
\(h_i=\ip{x}{\Omega_i}\le\sqrt{D_i}-.61\).
The target-independent inventory and helper nonadjacency in
Appendix~\ref{app:Acentralwindow} leave only pure far helpers or
near-four roles at \(5,6\) and their reflections.
Its older energy-dependent near-four conclusion is not used at the
larger budget \(.91\).

A left near-four has
\(v=\Omega_i/\sqrt2+\Omega_{i-1}/2\), \(i=5,6\).
The same actual wall is the opposite in direction \(j=i-1\).
Keeping the other original roots and all of \(Z\), positive finite
projection gives its actual endpoint, with
\[
 \Delta_j=1/(17-j),\qquad
 C_j=\sqrt{2j/17}\{(17-j)/2+(16-j)/\sqrt2\}.
\]
For \(j=4,5\), their squares are respectively greater than
\(177608/1700\) and \(187690/1700\), both greater than \(100\).
Thus the normalized mate height exceeds \(\sqrt{9/11}>.9\).
At the opposite end, either a second mate has height \(>.9\), or
nonadjacency supplies a strong \(11/12\) direction of height \(>.77\).
Their correlation is at most \(\sqrt{5/22}<1/2\).
Cauchy would give
\[
 1.67^2<e_A(2+2\rho)<3(.91)=2.73,
\]
whereas \(1.67^2=2.7889\).  Both orientations are impossible.
Every remaining weak site is therefore a nonadjacent pure helper.

Use the shared rational reference of Appendix~\ref{app:Aendpointcores}:
\[
 H_5=2.046,\ H_6=2.176,\ H_7=2.259,\ H_8=H_9=2.300,
 \qquad H_{17-i}=H_i.
\]
The ends are linear.  A weak end is flattened through six; every
other weak knot is removed by its neighboring chord before pressure.
The resulting positive-dual reference \(w\) has all coefficients at
actual strong floors and obeys
\[
\begin{array}{c|c}
 \text{number of weak ends among }5,12&E=\norm w^2\text{ lower}\\ \hline
 0&.86103645\\
 1&>.8285\\
 2&E_{\min}=1195276/1500000 .
\end{array}                                                   \tag{T3}
\]
The last rows include every permitted internal omission.
A weak end block has length \(L=6\), mean
\(b=H_6/(6\sqrt2)<.257\); a strong end block has \(L=5\),
mean \(b=H_5/(5\sqrt2)<.29\).  In both cases \(b>.25\),
the next coordinate is \(<.1\), and the individual strong-anchor
capacities are at least \(839/6000\) and \(349/2500\), respectively.

If an original signed endpoint coordinate were at least \(.565\),
replace its block by \(.565,r,\ldots,r\), keeping the block sum.
Here \(r=(Lb-.565)/(L-1)\) is greater than \(.17125\) or
\(.187\), so every surviving coefficient stays nonnegative.
The exact variance gain is \(L(.565-b)^2/(L-1)\).
For a weak end, its total with \(e_Z\) exceeds
\[
 E_{\min}+.09+(6/5)(.565-.257)^2>1.00068746.
\]
For a strong end it exceeds
\(.8285+.09+(5/4)(.565-.29)^2=1.01303125\).
Hence both signed endpoint coordinates are \(<.565\).
The actual edge condition in either direction is
\[
 x_{\rm end}\ge\sqrt{16/17}\sqrt{(C-1)/(C+1)}.
\]
At \(C=41/20\) its square is \(336/1037>.565^2\), with
positive cleared numerator \(198547\) over \(41480000\).
Therefore BOTH actual endpoint cosines satisfy
\[
                         C<41/20.                         \tag{T4}
\]
This is one reference contradiction per endpoint, not addition of
two independently optimized lower energies.

\subsubsection{The complete endpoint and foreign-contact list}
Release old \(1\), retaining the full old \(A_{15}\) and \(Z\).
Its actual opposite is isolated within \(A_{16}\), or has a single
three at \(j=2,3,15,16\), or a single four at \(j=2,16\).
Two contacts make a cycle; exceptional finite types have insufficient rank.
An old \(A_2\) admits at most one new contact, with complete costs
\[
                  Q=1/3,\quad2/3,\quad\phi^2/3.
\]
An old \(I_2(7)\) admits none, since even its smallest contact costs
\(1/[4\sin^2(\pi/7)]>1\).
When there is also an \(A_{15}\) contact, only an endpoint-three
with old \(A_2\) cost \(Q=1/3\) remains finite: the full chain
is \(A_{18}\).  A foreign arm at a \(D\) position has excessive
arms; a four or five becomes an internal multiple bond in the long path.

Write \(c=Ap-g\Omega_1-\beta\Omega_j-v_Z\), where
\(\beta=1/2\) or \(1/\sqrt2\).  For \(Q=0\), the exact data are
\[
\begin{array}{c|rrrr}
 \text{white three at }j&2&3&15&16\\ \hline
 C&(32g+15)/17&(16g+7)/\sqrt{17}&
   (16g+1)/\sqrt{17}&(32g+1)/17
\end{array}
\]
and white fours at \(2,16\) give
\((16g+15/\sqrt2)/\sqrt{17/2}\) and
\((16g+1/\sqrt2)/\sqrt{17/2}\).
The residuals are \(17/32,1/8,1/16\), for endpoint three,
\(D\)-position three and endpoint four.  The one mixed
endpoint-three multiplies its pure cosine by \(\sqrt{51/19}>8/5\).

At \(j=3\), own zero has the actual affine-\(E_8\) null
\[
 2n_1+4n_2+6n_3+5n_4+4n_5+3n_6+2n_7+n_8+3c;
\]
every positive own has \(C\ge15/\sqrt{17}>41/20\).
At \(j=15\), own zero has time square \(-2/17\), own three
has the null \(c+n_1+\cdots+n_{15}\), and owns at least four
violate (T4).  White four at \(2\) violates (T4) even at own zero;
at \(16\), own zero has time square \(-1/17\), and every
positive own violates (T4).

For pure white three at \(2\), own zero has time square \(-2/17\),
own three has the actual triangle null, and every own at least four
violates (T4).  The mixed version retains that triangle; its own-zero
null is the actual affine \(E_7\) combination
\[
 2n_1+4n_2+3n_3+2n_4+n_5+3c+2z_1+z_2,
 \qquad \ip c{z_1}=\ip{z_1}{z_2}=-1/2.
\]
Mixed white three at \(16\) has own-zero time square \(-10/51\),
the own-three all-three cycle, and for own at least four has
\(C>(8/5)(32(.7)+1)/17=936/425>41/20\).
A parallel own always has the null \(c+n_1\).
Every displayed null has positive cap time and nonpositive product
against every omitted wall.  A positive finite projection likewise
checks all unselected walls termwise.

The only remaining nonisolated endpoint is consequently PURE far three:
\[
 c=Ap-g\Omega_1-\Omega_{16}/2,\quad
 A^2=(32g^2+2g-9)/17,
 \quad 0<g<677/640<53/50,\quad A^2<7/4.              \tag{T5}
\]
Own zero has negative time, own three has zero time, and reflection
gives the other end.  Finite and dotted own values remain present.

\subsubsection{Using the actual defect wall}
For \(N_7\), a pure endpoint wall is distinct from \(c_Z\) and
has actual product \(-A\sqrt{h_N}\) with it.  The native gap
forces \(A^2h_N\ge1/4\), so \(A^2>3\), excluding (T5).
An isolated-within-A endpoint instead has
\[
 A^2=Dg^2+Q-1,\quad C^2=Dg^2/(1-Q),\quad
 D=32/17,\quad u=(1-Q)/g^2=D/C^2.
\]
Own zero has negative time, since the full white cost satisfies \(Q<1\).
At \(Q=0\), the same gap gives \(u<D/4<1\).
At \(Q>0\), a contact \(\beta\) to the old \(A_2\) gives
\(Q=4\beta^2/3\) and actual defect cross
\(q=4g_7\beta/3\) or \(2g_7\beta/3\).
Acuteness, including orthogonality, gives
\[
 A^2\ge q^2/h_N\ge Q(2+\theta)/[4(\theta-1)]>13Q/4.
\]
Since \(Q\ge1/3\),
\(C^2>[1+(9/4)Q]/(1-Q)\ge21/8>D\), so again \(0<u<1\).
This covers every foreign orientation and finite or dotted own.

For \(O_7\) all endpoints are pure.  Set \(M=4A^2/k\).
It is four times the square of an ACTUAL negative two-wall product.
Its positive values therefore begin
\(1,2,\phi^2,3,2+\theta\), followed by larger finite values,
parallel \(4\) or dotted values greater than four.
For the far helper, own four gives \(M<1\); finite owns at least
five give \(1<M<2\).  In the dotted case, (T5) gives
\(1<M<7/3<\phi^2\), so the only possibility is \(M=2\).
It forces
\[
 A^2=k/2,\quad 32g^2+2g-9=17k/2,\quad
 1.013<g<1.014,\quad1.2346<A<1.235.                 \tag{T6}
\]
These brackets follow by positive substitution using (T2).

Two far helpers would have this same \(g\).  They are distinct because
their own \(g>1\) differs from the other-end helper coefficient \(1/2\).
Their spatial
product \((2g^2+32g+1/2)/17>69/34\) exceeds their time product
\(A^2=k/2<61/40\).  This violates acuteness.  Thus a far helper
would face an isolated endpoint.
If that isolated endpoint had finite own \(\delta\le1\), then
\(0<M_\delta=4(D\delta^2-1)/k<20/17<2\), so it must equal one.
For labels at most twelve,
\(4(D\delta^2-1)\le(32\sqrt3-4)/17<3.0485<k\).
For labels at least thirteen, \(\delta>.97\) gives
\(4(D\delta^2-1)>3.0844>k\).  No finite label works; own zero
has negative time and parallel is the actual null pair.
The other endpoint is therefore dotted, with \(\delta>1\) and
\(\sqrt D\delta<41/20\).

Put \(B=(2g+16)/17\in(1.06,1.0605)\).
Their actual mutual magnitude
\[
 m(\delta)=A\sqrt{D\delta^2-1}-B\delta
\]
increases on \(\delta\ge1\), because its derivative exceeds
\(1.2346(1.37)-1.0605>0\).
It is positive at one, and its endpoint upper is
\(1.235(1.79)-1.06(1.49)=.63125<.7\).
The native gap thus forces \(m=1/2\).
At \(1.35\) it is \(<.4956\), while at \(1.358\) it is
\(>.5006322\); hence \(1.35<\delta<1.358\).
But its own ACTUAL defect-wall product now has
\[
 3<M_\delta<(1.313)(2.472)=3.245736<2+\theta,
\]
the forbidden native-six/native-seven gap.  This excludes the last
far helper.  Both \(O_7\) endpoints are isolated and dotted.

\subsubsection{One simultaneous endpoint payment}
For either isolated \(N_7\) endpoint clear its full foreign vector
with its positive old-\(A_2\) inverse.  For \(O_7\), put \(Q=0\).
Each end has its OWN parameter \(u\in(0,1)\), and the actual core
\[
 \eta=c+v_Z+gu n_1
 =g\{\sqrt{D-u}\,p-Z(u)\},\qquad Z(u)=\Omega_1-u n_1
\]
has square \(-g^2u(1-u)\), positive time, zero selected cap and
foreign products, and negative \(n_1\) product.  All other walls,
including the opposite endpoint and the defect wall, are nonpositive
termwise.  Its actual normalization gives
\[
 \ip x{Z(u)}\ge h(u)=\sqrt{D-u}-\sqrt{u(1-u)}.
\]
At a length-\(L\) block put \(Y=Z-\Omega_L/L\), with the
reflected signed expression at the right end.  Then
\[
 \ip wY=0,\quad \ip{Y_l}{Y_r}=0,\quad
 \ip wZ=H_L/L\le.4092,\quad
 Y^2=2-2/L-2u+u^2\le5/3-2u+u^2.
\]
The residuals have zero sum and disjoint first/last block supports.
Pay \(\mu=.32\) separately at BOTH surviving strong anchors:
\(\mu/L\le.064\) is below both capacities in (T3).
All old coefficients remain nonnegative before either pressure is used.
For \(W=w+.32(Y_l+Y_r)\), one original-vector square yields
\[
 e_A>E+G(u_l)+G(u_r),\quad
 G(u)=.64(h(u)-.4092)-.1024(5/3-2u+u^2).            \tag{T7}
\]
The spectator energy has not been added to either gain.

For completeness the following finite scalar certificate covers every
independent \(0\le u\le1\).  Write
\[
 G_0=.64(h(u)-.4112)-.1024(8/5-2u+u^2),\qquad
 G=G_0-208/37500.
\]
Positive square tests for \(\sqrt{D-u}\) at \(0,1/2,3/4,1\)
give lower bounds \(1.3719,1.1757,1.0641,.9393\).
Use the root chords and \(u^2\le(l+r)u-lr\) on each interval:
\[
\begin{array}{c|c|c|c|c}
[l,r]&\text{root chord}&K&a&\sqrt{a^2+.64^2}\text{ upper}\\ \hline
[0,1/2]&1.3719-.3924u&.451008&.097536&.6475\\
[1/2,3/4]&1.3989-.4464u&.506688&.208896&.6733\\
[3/4,1]&1.4385-.4992u&.570432&.293888&.7043
\end{array}
\]
Here \(G_0>K-au-.64\sqrt{u(1-u)}\).
The four root squares are
\(1.88210961,1.38227049,1.13230881,.88228449\),
each below \(32/17,47/34,77/68,15/17\), respectively.
The three upper-radical square gaps are
\(.000142978704,.000095351184,.000068333456\), all positive.
Finally
\(au+.64\sqrt{u(1-u)}\le(a+\sqrt{a^2+.64^2})/2\),
so the three lower bounds for \(G_0\) are
\(.07849,.06559,.071338\).  Thus, uniformly,
\[
               G>.06559-208/37500=18013/300000>.06.
\]
Combining (T1), (T3) and (T7) now gives
\[
 e_A+e_Z>E_{\min}+.12+.09>1.00685066>1,
\]
a contradiction.  The two end parameters were independent, all actual
walls were retained, and the rank-two energy was counted once.
This excludes BOTH actual typed-seven alternatives beside \(A_{16}\).

Together with the already proved uniform \(A_{16}>.86\) and the
ordinary/joint spectator bounds, this closes its entire exact-eighteen
row: ordinary \(T_2\), two ordinary \(A_1\)'s and \(I+A_1\)
are excluded by those earlier bounds, and the two typed cases by the
argument here.  It does not assert an arbitrary-spectator
\(A_{16}>.91\) theorem.
